% 经典 Beamer 模板：Madrid 内置主题 + 默认配色，只用标准 pdflatex 和 TeX Live 自带的包。
% 复制整个 template/ 目录后，改标题区和 frame 内容即可；版式、页脚、导航条都不要动。
\documentclass[aspectratio=169,10pt]{beamer}
\usetheme{Madrid}
\usecolortheme{default}

\usepackage{booktabs}
\usepackage{multirow}
\usepackage{array}
\usepackage{graphicx}
\usepackage{amssymb}

% ---- 常用宏 -------------------------------------------------------
\newcommand{\tbd}{\textcolor{gray}{TBD}}                            % 还没出的结果
\newcommand{\sd}[1]{{\scriptsize\textcolor{gray}{$\pm$#1}}}         % 标准差
\newcommand{\ok}{\textcolor{green!50!black}{done}}
\newcommand{\pend}{\textcolor{orange!80!black}{pending}}
\newenvironment{takeaway}{\begin{exampleblock}{Takeaway}}{\end{exampleblock}}   % 只放在关键页
\newcolumntype{L}[1]{>{\raggedright\arraybackslash}p{#1}}           % 左对齐定宽列（长文字表格用）

% ---- 标题区：副标题只放论文 / 报告副标题，汇报性质写在日期行 -----------
\title[Short Title]{Full Title of the Talk}
\subtitle{One-Line Subtitle}
\author[A.~Author]{Author Name}
\institute[Univ.]{Institution}
\date[Jan 1, 2027]{Progress Report\\[4pt] January 1, 2027}

\begin{document}

\begin{frame}
  \titlepage
\end{frame}

% Outline：两列，每个 section 下列出 subsection（每个 frame 前放一个 \subsection）
{\addtobeamertemplate{section in toc}{\vspace{4pt}}{}
\begin{frame}{Outline}
  \begin{columns}[T]
    \begin{column}{0.48\textwidth}
      \tableofcontents[sections={1-2}]
    \end{column}
    \begin{column}{0.48\textwidth}
      \tableofcontents[sections={3-4}]
    \end{column}
  \end{columns}
\end{frame}
}

% ------------------------------------------------------------------
\section{Motivation}

\subsection{Problem}
\begin{frame}{Problem Statement}
  \begin{itemize}
    \item One idea per bullet, three to five bullets per slide
    \item Explain a technical term in plain words the first time it appears
    \item Titles are statements or plain topic names, not questions
  \end{itemize}
  \begin{takeaway}
    One sentence the audience should remember from this section.
  \end{takeaway}
\end{frame}

\subsection{Overview}
\begin{frame}{Contributions and Evidence}
  \centering\small
  \begin{tabular}{@{}L{0.32\textwidth}L{0.44\textwidth}L{0.17\textwidth}@{}}
    \toprule
    Contribution & Evidence & Status \\
    \midrule
    \addlinespace (1) First claim & experiment or analysis that supports it & \ok \\
    \addlinespace (2) Second claim & experiment or analysis that supports it & \pend{} (date) \\
    \bottomrule
  \end{tabular}
\end{frame}

% ------------------------------------------------------------------
\section{Results}

\subsection{Main table}
\begin{frame}{Main Results}
  \centering\small
  \begin{tabular}{@{}lcccc@{}}
    \toprule
    & \multicolumn{2}{c}{Dataset A} & \multicolumn{2}{c}{Dataset B} \\
    \cmidrule(lr){2-3}\cmidrule(lr){4-5}
    Model & metric 1 & metric 2 & metric 1 & metric 2 \\
    \midrule
    Baseline & 35.5\sd{0.8} & 84.1 & 24.0\sd{0.7} & 83.8 \\
    Ours & \textbf{37.7}\sd{0.8} & \textbf{76.6} & \tbd & \tbd \\
    \bottomrule
  \end{tabular}\\[2pt]
  {\scriptsize Footnote: what the metric is, what $\pm$ means, how many questions.}
  \begin{takeaway}
    The single comparison that matters, with its number.
  \end{takeaway}
\end{frame}

\subsection{Figure}
\begin{frame}{Figure with Explanation}
  \begin{columns}[T]
    \begin{column}{0.55\textwidth}
      \centering
      % \includegraphics[width=\linewidth]{fig/figure.pdf}
      \fbox{\parbox[c][0.5\textheight][c]{0.9\linewidth}{\centering figure: fig/figure.pdf}}
    \end{column}
    \begin{column}{0.43\textwidth}
      \small
      \begin{itemize}
        \item What the axes or panels show
        \item The one pattern to notice
      \end{itemize}
    \end{column}
  \end{columns}
\end{frame}

% ------------------------------------------------------------------
\section{Plan}

\subsection{Timeline}
\begin{frame}{Remaining Work and Timeline}
  \begin{columns}[T]
    \begin{column}{0.36\textwidth}
      \small
      \textbf{Work}
      \begin{itemize}
        \item First open item
        \item Second open item
      \end{itemize}
    \end{column}
    \begin{column}{0.62\textwidth}
      \small
      \textbf{Timeline} (JST)\\[2pt]
      \footnotesize
      \begin{tabular}{@{}ll@{}}
        \toprule
        Date & Item \\
        \midrule
        Jan 1--3 & one line per period \\
        Jan 10 20:59 & \textbf{deadline} \\
        \bottomrule
      \end{tabular}
    \end{column}
  \end{columns}
\end{frame}

% ------------------------------------------------------------------
\section{Summary}

\subsection{Summary}
\begin{frame}{Summary}
  \begin{enumerate}
    \item First finding, with its key number
    \item Second finding
    \item What is still open
  \end{enumerate}
\end{frame}

\end{document}
